\begin{defi}\label{def:quiver}
Given a triangulation $\mathcal{T}$ of $C(n + 2d + 1, 2d)$, we define the \emph{quiver} $Q(\mathcal{T})$ of $\mathcal{T}$ to be the following directed graph. The vertices are \[Q_{0}(\mathcal{T}) = \darcs{\mathcal{T}}.\] The arrows are given by $A \to B$ with $A \neq B$ such that $(A - \onevec) \swr B$ and there exists no $A' \in \darcs{\mathcal{T}}\setminus\{A, B\}$ with $(A - \onevec) \swr A'$ and $(A' - \onevec) \swr B$. Recall our convention of taking $A - \onevec$ modulo $n + 2d + 1$.
\end{defi}

We define the quiver in this way so that it coincides with the Gabriel quiver of the endmorphism algebra of the cluster-tilting object corresponding to the triangulation. The arrows mirror the description of the homomorphisms in the cluster category $\mathcal{O}_{A_{n}^{d}}$ from \cite{ot}. However, we now prove a simpler description of the quiver, given in Corollary~\ref{cor:quiv_desc}. In order to obtain this, we first make some observations about cyclic polytopes. Changing the ordering on on $[n + 2d + 1]$ from $<_{1}$ to $<_{l}$ induces a (combinatorial) automorphism of the cyclic polytope $C(n + 2d + 1, 2d)$; see \cite{kw}. We think of this automorphism as reorienting the cyclic polytope. Generalising \cite{ot}, given a $2d$-simplex $T$ of $C(n + 2d + 1,2d)$ and an ordering $<_{l}$ of $[n + 2d + 1]$ under which $T$ is ordered $T = (t_{0}, t_{1}, \dots , t_{2d})$, we write $e_{l}(T) = (t_{0}, t_{2}, \dots , t_{2d-2}, t_{2d})$. If the ordering is determined by the context, we simply write $e(T)$.

\begin{lemma}\label{lem:find_2d_simplex}
Given an ordering $<_{l}$ of $[n + 2d + 1]$ and a $d$-arc $A$ of a triangulation $\mathcal{T}$ of $C(n + 2d + 1,2d)$, there is a $2d$-simplex $T$ of $\mathcal{T}$ such that $e_{l}(T) = A$.
\end{lemma}
\begin{proof}
This follows from applying \cite[Proposition 2.13]{ot} in the orientation given by~$<_{l}$.
\end{proof}

\begin{prop}\label{prop:all_but_one_vertex}
Suppose that $A \rightarrow B$ is an arrow in $Q(\mathcal{T})$. Then $A$ and $B$ share all but one entry, and there is a $2d$-simplex $T$ of $\mathcal{T}$ such that $A$ and $B$ are both faces of $T$.
\end{prop}
\begin{proof}
We have that $(A - \onevec) \swr B$. It must be the case that $A$ and $B$ have at least one common entry, otherwise $A \swr B$. Hence, by re-ordering, we may assume that $a_{0} = b_{0}$ and $a_{d - 1} < a_{d} < b_{d}$. There must be a $2d$-simplex $T$ of $\mathcal{T}$ such that $e(T) = B$. We label the vertices of $T$ by $T = (b_{0}, x_{1}, b_{1}, x_{2}, \dots, b_{d - 1}, x_{d}, b_{d})$.

\sloppy We cannot have that $b_{i - 1} < x_{i} \leq a_{i} - 1$ for all $i \in [d]$, otherwise $A \swr (x_{1}, x_{2}, \dots, x_{d}, b_{d})$, which is impossible since $A$ and $(x_{1}, x_{2}, \dots, x_{d}, b_{d})$ are both $d$-arcs of $\mathcal{T}$. Hence there is an $i \in [d]$ such that $b_{i - 1} < a_{i} - 1 < x_{i}$. Then $(b_{0}, b_{1}, \dots, b_{i - 1}, x_{i}, b_{i + 1}, b_{i + 2}, \dots, b_{d})$ is a $d$-arc of $\mathcal{T}$, since $b_{i - 1} < a_{i} - 1 < x_{i}$ and $x_{i} < b_{i} < b_{i + 1}$. Moreover, $(A - \onevec) \swr (b_{0}, b_{1}, \dots, b_{i - 1}, x_{i}, b_{i + 1}, b_{i + 2}, \dots, b_{d})$ and $(b_{0} - 1, \dots, b_{i - 1} - 1, x_{i} - 1, b_{i + 1} - 1, \dots, b_{d} - 1) \swr B$. Since $A \rightarrow B$ is an arrow in $Q(\mathcal{T})$, we must have $A = (b_{0}, b_{1}, \dots, b_{i - 1}, x_{i}, b_{i + 1}, b_{i + 2}, \dots, b_{d})$. (In fact, by the ordering we have chosen, we must have that $i = d$.) Therefore $A$ and $B$ are both faces of the $2d$-simplex $T$ and they share all but one entry, as desired.
\end{proof}

\begin{coro}\label{cor:quiv_desc}
The arrows of $Q(\mathcal{T})$ are \[Q_{1}(\mathcal{T}) = \set{A \rightarrow A + \rvec \st \parbox{6.5cm}{\begin{center}$A, A + \rvec \in \nonconsec{n + 2d + 1}{d},$ $\nexists s \in [r - 1] \text{ such that } A + \svec \in Q_{0}(\mathcal{T})$\end{center}}}.\]
\end{coro}

We call arrows $A \to A + \rvec$ \emph{arrows of type $i$}, analogously to our terminology of arrows of type~$i$ in $\qdn$ from Section~\ref{sect:back:cut_slice}. Note that this labelling depends on the choice of an ordering $<_{l}$.

For an arrow $\alpha$, we denote the \emph{head} $h(\alpha)$ and the \emph{tail} $t(\alpha)$ such that $t(\alpha) \xrightarrow{\alpha} h(\alpha)$. We say that $\alpha$ is \emph{incident} at $t(\alpha)$ and $h(\alpha)$. Given a quiver $Q$ with vertices $A,B \in Q_{0}$, by a \emph{path} in $Q$ from $A$ to $B$ we mean a finite sequence of arrows $\alpha_{1}\dots\alpha_{r}$ such that $t(\alpha_{1}) = A, h(\alpha_{r}) = B$ and $h(\alpha_{i - 1}) = t(\alpha_{i})$ for all $i \in \{2, 3, \dots, s\}$. If there is a path from $A$ to $B$, then we write $A \leadsto B$. We note the following property concerning paths in $Q(\mathcal{T})$ which will be useful later.

\begin{lemma}\label{lem:find_path}
Given a triangulation $\mathcal{T}$ of $C(n + 2d + 1, 2d)$ and $A, B \in Q_{0}(\mathcal{T})$ with $(A - \onevec) \swr B$, there is a path from $A$ to $B$ in $Q(\mathcal{T})$.
\end{lemma}
\begin{proof}
This is clear from Definition~\ref{def:quiver}, using induction.
\end{proof}
